%% 第 14 章：实向量空间的复化
%% 本文件由 tools/md2tex.py 自动生成，请勿直接编辑。

\chapter{实向量空间的复化}


\section{为什么需要复化}

复数域上的有限维算子总能上三角化，并拥有 Jordan 标准型。

实数域上则不一定有实特征值。

例如，平面旋转算子没有非零实特征向量。因此，它无法在实数基下被对角化，也无法在实数基下上三角化。

复化的作用是：把实空间扩展到复数域，使原本隐藏的复特征值与复 Jordan 链显现出来。

\begin{jdefn}{复化空间}{r:15:1}
设 $V$ 是实向量空间。

定义其复化：

\[
V_{\mathbb C}
=
\{x+iy:x,y\in V\}.
\]

它是一个复向量空间。

若：

\[
T:V\to V
\]

是实线性算子，则定义其复化算子：

\[
T_{\mathbb C}:V_{\mathbb C}\to V_{\mathbb C},
\]

\[
T_{\mathbb C}(x+iy)
=
T(x)+iT(y).
\]
\end{jdefn}

\section{复化的矩阵意义}

若 $T$ 在实基下的矩阵是实矩阵 $A$，则 $T_{\mathbb C}$ 在对应复基下的矩阵仍然是同一个 $A$。

变化不在矩阵本身，而在于允许输入、输出坐标取复数。

\begin{jprop}{共轭特征值}{r:15:2}
由于 $T$ 的矩阵系数都是实数，若：

\[
\lambda=a+ib
\]

是 $T_{\mathbb C}$ 的特征值，则：

\[
\overline\lambda=a-ib
\]

也是特征值。

相应的 Jordan 结构也成对出现。
\end{jprop}
